\documentclass[11pt]{article}
\usepackage[T1]{fontenc}
\usepackage{lmodern}
\usepackage{amsmath,amssymb,amsthm,mathtools}
\usepackage[margin=1in]{geometry}
\usepackage[hidelinks]{hyperref}

\newtheorem{theorem}{Theorem}
\newtheorem{proposition}{Proposition}
\newtheorem{corollary}{Corollary}
\newtheorem{remark}{Remark}
\DeclareMathOperator*{\esssup}{ess\,sup}
\newcommand{\C}{\mathbb{C}}
\newcommand{\R}{\mathbb{R}}
\newcommand{\eps}{\varepsilon}
\newcommand{\Reps}{R_{\eps}}
\newcommand{\Deps}{D_{\eps}}
\newcommand{\Thetaeps}{\Theta_{\eps}}
\newcommand{\aPlus}{\alpha_{+}}
\newcommand{\aMinus}{\alpha_{-}}

\title{M2 + E17: Positivity and Singular-Limit Consistency of the
Regularized Curvature Operator}
\author{}
\date{}

\begin{document}
\maketitle

\begin{abstract}
We prove that the corrected modulus-squared regularization of the WCT
curvature-feedback operator is well-defined for all finite complex field
values at every fixed regularization scale $\eps>0$. The regularized
reciprocal is smooth and globally bounded, converges to the ordinary
reciprocal at every nonzero field value, and admits quantitative uniform
and $L^2$ convergence estimates away from field nodes. We also verify
global $U(1)$ phase invariance and exhibit a counterexample to the
historical, uncorrected denominator. These algebraic and analytic
properties do not by themselves imply well-posedness or stability of a
nonlinear evolution equation.
\end{abstract}

\section{Definitions and theorem}

Let $\Omega\subseteq\R^d$ be open, $\alpha\in\R$, and $\eps>0$. For
$z\in\C$, define
\begin{equation}\label{eq:definitions}
\Deps(z):=|z|^2+\eps^2e^{-2\alpha|z|^2},
\qquad
\Reps(z):=\frac{\overline z}{\Deps(z)}.
\end{equation}
For $\psi\in C^2(\Omega;\C)$, define the regularized curvature operator
\begin{equation}\label{eq:curvature}
\Thetaeps[\psi](x):=-(\Delta\psi(x))\Reps(\psi(x)).
\end{equation}
The Laplacian is applied componentwise to the complex-valued field.

\begin{theorem}[M2 + E17]\label{thm:main}
For every $\eps>0$ and $\alpha\in\R$ the following hold.
\begin{enumerate}
\item $\Deps(z)>0$ for every $z\in\C$, including $z=0$.
\item $\Reps$ is a globally bounded $C^\infty$ map $\R^2\to\R^2$,
identifying $\C$ with $\R^2$, and $\Reps(0)=0$.
\item For every $z\ne0$,
\begin{equation}\label{eq:limit}
\lim_{\eps\to0^+}\Reps(z)=\frac{1}{z}.
\end{equation}
\item If $\alpha\ge0$, then for every $\delta>0$,
\begin{equation}\label{eq:uniform}
\sup_{|z|\ge\delta}
\left|\Reps(z)-\frac1z\right|
\le\frac{\eps^2}{\delta^3}.
\end{equation}
For arbitrary $\alpha\in\R$, the convergence is uniform on every
compact subset of $\C\setminus\{0\}$.
\item For $\psi\in C^2(\Omega;\C)$ and every $x\in\Omega$ such that
$\psi(x)\ne0$,
\begin{equation}\label{eq:curvaturelimit}
\lim_{\eps\to0^+}\Thetaeps[\psi](x)
=-\frac{\Delta\psi(x)}{\psi(x)}.
\end{equation}
\end{enumerate}
\end{theorem}

\begin{proof}
\emph{Step 1: Strict denominator positivity.}
For $z\in\C$, $|z|^2\ge0$. Because $\alpha$ and $|z|^2$ are real,
$e^{-2\alpha|z|^2}>0$. Multiplying by $\eps^2>0$ gives
\[
\Deps(z)=\underbrace{|z|^2}_{\ge0}
+\underbrace{\eps^2e^{-2\alpha|z|^2}}_{>0}>0.
\]
At the node $z=0$, $\Deps(0)=\eps^2>0$, and consequently
$\Reps(0)=0$.

\emph{Step 2: Smoothness and boundedness.}
Writing $z=x+iy$, the numerator $x-iy$ and the denominator
\[
x^2+y^2+\eps^2e^{-2\alpha(x^2+y^2)}
\]
are $C^\infty$ in the real variables $(x,y)$. Step 1 shows that the
denominator is nowhere zero, so $\Reps\in C^\infty(\R^2;\R^2)$.

To prove a global bound, fix any $\rho>0$ and set
$\aPlus:=\max\{\alpha,0\}$. When $|z|\le\rho$,
\[
\Deps(z)\ge \eps^2e^{-2\aPlus\rho^2},
\qquad
|\Reps(z)|\le
\frac{\rho e^{2\aPlus\rho^2}}{\eps^2}.
\]
When $|z|\ge\rho$, $\Deps(z)\ge|z|^2$, hence
\[
|\Reps(z)|\le\frac{|z|}{|z|^2}
=\frac1{|z|}\le\frac1\rho.
\]
Therefore
\begin{equation}\label{eq:globalbound}
\|\Reps\|_{L^\infty(\C)}
\le\max\!\left\{\frac1\rho,
\frac{\rho e^{2\aPlus\rho^2}}{\eps^2}\right\}<\infty.
\end{equation}
This bound is for each fixed $\eps>0$; it is not asserted to be
uniform as $\eps\to0^+$.

\emph{Step 3: Reciprocal limit.}
For fixed $z\ne0$, the quantity $e^{-2\alpha|z|^2}$ is finite and
independent of $\eps$. Thus
\[
\lim_{\eps\to0^+}\Deps(z)=|z|^2>0.
\]
Continuity of division at a nonzero denominator yields
\[
\lim_{\eps\to0^+}\Reps(z)
=\frac{\overline z}{|z|^2}
=\frac{\overline z}{z\overline z}
=\frac1z.
\]

\emph{Step 4: Quantitative convergence away from nodes.}
For $z\ne0$, direct subtraction gives the exact identity
\begin{align}
\Reps(z)-\frac1z
&=\frac1z\left(
\frac{|z|^2}{|z|^2+\eps^2e^{-2\alpha|z|^2}}-1\right)
\notag\\
&=-\frac1z\,
\frac{\eps^2e^{-2\alpha|z|^2}}
{|z|^2+\eps^2e^{-2\alpha|z|^2}}.
\label{eq:exactdifference}
\end{align}
Therefore
\begin{equation}\label{eq:error}
\left|\Reps(z)-\frac1z\right|
=\frac{\eps^2e^{-2\alpha|z|^2}}
{|z|\bigl(|z|^2+\eps^2e^{-2\alpha|z|^2}\bigr)}
\le\frac{\eps^2e^{-2\alpha|z|^2}}{|z|^3}.
\end{equation}
If $\alpha\ge0$ and $|z|\ge\delta>0$, then the exponential is at
most one. This proves \eqref{eq:uniform}.

For general $\alpha$, let $K\subset\C\setminus\{0\}$ be compact.
There exist $0<\delta\le M<\infty$ such that
$\delta\le|z|\le M$ on $K$. With
$\aMinus:=\max\{-\alpha,0\}$, we have
\[
\sup_{z\in K}
\left|\Reps(z)-\frac1z\right|
\le\frac{\eps^2e^{2\aMinus M^2}}{\delta^3}
\xrightarrow[\eps\to0^+]{}0.
\]

\emph{Step 5: Curvature limit.}
Fix $x\in\Omega$ with $\psi(x)\ne0$. Since $\psi$ and
$\Delta\psi$ are fixed independently of $\eps$, Step 3 gives
\[
\lim_{\eps\to0^+}\Thetaeps[\psi](x)
=-(\Delta\psi(x))\lim_{\eps\to0^+}\Reps(\psi(x))
=-\frac{\Delta\psi(x)}{\psi(x)}.
\]
All the assertions follow.
\end{proof}

\section{Convergence in $L^2$}

\begin{corollary}[Strong $L^2$ limit away from nodes]
\label{cor:l2}
Assume $\alpha\ge0$,
$\psi\in H^2(\Omega;\C)$, and
$|\psi(x)|\ge\delta>0$ for almost every $x\in\Omega$.
Then
\begin{equation}\label{eq:l2}
\left\|\Thetaeps[\psi]
+\frac{\Delta\psi}{\psi}\right\|_{L^2(\Omega)}
\le\frac{\eps^2}{\delta^3}
\|\Delta\psi\|_{L^2(\Omega)}
\xrightarrow[\eps\to0^+]{}0.
\end{equation}
\end{corollary}

\begin{proof}
By definition,
\[
\Thetaeps[\psi]+\frac{\Delta\psi}{\psi}
=(\Delta\psi)\left(\frac1\psi-\Reps(\psi)\right)
\quad\text{a.e. on }\Omega.
\]
Apply the pointwise multiplier estimate \eqref{eq:uniform} with
$z=\psi(x)$, then take the $L^2$ norm.
\end{proof}

\begin{corollary}[Fixed-regularization $L^2$ control]
If $\psi\in H^2(\Omega;\C)$ and $\eps>0$ is fixed, then
\[
\|\Thetaeps[\psi]\|_{L^2(\Omega)}
\le\|\Reps\|_{L^\infty(\C)}
\|\Delta\psi\|_{L^2(\Omega)}<\infty.
\]
\end{corollary}

\begin{proof}
The multiplier $\Reps(\psi)$ is essentially bounded by
\eqref{eq:globalbound}, while $\Delta\psi\in L^2(\Omega)$.
\end{proof}

\begin{remark}[Whole-space limitation]
The hypothesis $|\psi|\ge\delta>0$ almost everywhere on
$\Omega=\R^d$ is incompatible with $\psi\in L^2(\R^d)$, since it
would imply $\|\psi\|_2^2\ge\delta^2|\R^d|=\infty$.
Thus Corollary~\ref{cor:l2} is not a blanket global convergence
theorem for positive-mass localized whole-space fields.
\end{remark}

\begin{remark}[No uniform convergence through a node]
At $z=0$, $\Reps(0)=0$ whereas $1/z$ is undefined. The obstruction
is quantitative: set $z=\eps$ (real and positive). From
\eqref{eq:error},
\[
\left|\Reps(\eps)-\frac1\eps\right|
=\frac{e^{-2\alpha\eps^2}}
{\eps\bigl(1+e^{-2\alpha\eps^2}\bigr)}
\sim\frac{1}{2\eps}\quad(\eps\to0^+).
\]
Therefore the reciprocal convergence is not uniform in any punctured
neighborhood of the origin.
\end{remark}

\section{Global phase invariance}

\begin{proposition}[Constant $U(1)$ phase invariance]
For every constant $\beta\in\R$ and
$\psi\in C^2(\Omega;\C)$,
\[
\Thetaeps[e^{i\beta}\psi]=\Thetaeps[\psi].
\]
\end{proposition}

\begin{proof}
The denominator depends only on the modulus, so
\[
\Reps(e^{i\beta}z)=e^{-i\beta}\Reps(z).
\]
Since $\beta$ is constant,
$\Delta(e^{i\beta}\psi)=e^{i\beta}\Delta\psi$. Hence
\[
\Thetaeps[e^{i\beta}\psi]
=-e^{i\beta}(\Delta\psi)e^{-i\beta}\Reps(\psi)
=\Thetaeps[\psi].
\]
\end{proof}

\begin{remark}
Global $U(1)$ invariance does not by itself establish local gauge
invariance under a position-dependent phase $\beta(x)$.
\end{remark}

\section{Counterexample to the historical denominator}

The earlier expression
\begin{equation}\label{eq:old}
\Thetaeps^{\mathrm{old}}[\psi]
=-\frac{\Delta\psi}
{\psi+\eps e^{-\alpha|\psi|^2}}
\end{equation}
does not have a universally nonzero denominator. Indeed, take
$\alpha=0$ and $\psi\equiv-\eps$. Then
\[
\psi+\eps e^{-\alpha|\psi|^2}
=-\eps+\eps=0.
\]
The expression \eqref{eq:old} is undefined there, even if
$\Delta\psi=0$ (which gives the indeterminate quotient $0/0$).
By contrast, the corrected denominator satisfies
\[
|\psi|^2+\eps^2e^{-2\alpha|\psi|^2}>0
\quad\text{at every finite complex field value.}
\]
The historical denominator also fails generic constant-phase
covariance because its two summands transform differently.

\section{Scope of the result}

The theorem establishes denominator positivity, pointwise regularity
for fixed $\eps>0$, reciprocal consistency, quantitative convergence
away from nodes, and global phase invariance. It does \emph{not}
by itself prove nonlinear PDE well-posedness, global-in-time existence,
uniqueness, confinement, Lyapunov descent, or dynamical stability.
Those require separate hypotheses and proofs for the chosen evolution.

\begin{thebibliography}{9}
\bibitem{corrected}
R.~J. Reyes, \emph{WCT Full Equation List (Corrected)},
\url{https://github.com/rickyjreyes/geometry_of_resonance/blob/main/WCT_FULL_EQUATION_LIST_CORRECTED.md}.
\bibitem{master}
R.~J. Reyes, \emph{WCT Master Equations (Updated)},
\url{https://github.com/rickyjreyes/geometry_of_resonance/blob/main/WCT_MASTER_EQUATIONS_UPDATED.md}.
\bibitem{closure}
R.~J. Reyes, \emph{Derivation Closure},
\url{https://github.com/rickyjreyes/geometry_of_resonance/blob/main/DERIVATION_CLOSURE.md}.
\end{thebibliography}

\end{document}
